\begin{lema}[Regularity of the profile at the endpoint]\label{lem:regularidad-extremo}
For some \(x_0>0\), \(G\) is analytic in a neighborhood of \([0,x_0]\). In this neighborhood, \(\numerr\)
is meromorphic with its only possible pole at zero, of order at most four.
\end{lema}
\begin{proof}
The continuous unit-mass profile satisfies \(\nu_1\ge c>0\), has unit axial marginals, and satisfies
\(u_f\ge c>0\) in both oblique families~\cite{pv}. Reduce \(c\) so that \(c<1\). On a complex disk
\(|x|<r\) we have \(|U_f|\ge c/2\), so \(W\) and \(1/U_f\) vary holomorphically on this disk as bounded
functions of position. All the domains of crossings, fibers, parities, and incidences in
Section~\ref{sec:contracciones} are independent of \(x\).

\paragraph{The inverse and its entries.}
Use a fixed measure \(d\sigma=dy\) on each axis and \(dy/2\) on each oblique sheet. Set \(D_x(a)=1\) on
the axes and \(D_x(a)=x\) on the obliques, and let \(u_x(a)=1\) or \(U_a(x)\), respectively. Thus
\(d\rho=(u_x/D_x)d\sigma\). For two distinct intersecting lines, \(K(a,b)=\nu_1(a\cap b)\eta_a\eta_b\),
with \(\eta_a=D_x(a)/u_x(a)\). Hence \(K=D_x K_0 D_x\), where the kernel \(K_0\) is holomorphic and
uniformly bounded, including at the edges under the fixed crossing convention. The remaining entries are
zero. In \(L^2(d\sigma)\), kernel products use \(d\sigma\), and the factors \(D_x,u_x\) act by
multiplication. The operator to be inverted is
\[
 \mathcal T_x=I+u_x K_0D_x+\Pi .
\]
At \(x=0\) it is block triangular (axes first), with the identity on the oblique block. The axial block
is the identity plus the operator with cross blocks \(\nu_1,\nu_1^{\mathsf T}\) plus \(\Pi\). On the
common and alternating axial constants its eigenvalues are \(2,1\). On functions with mean zero on each
axis the norm of the cross-block operator is \(\le1-c\). The kernel \((\nu_1-c)/(1-c)\) is doubly
stochastic, hence a contraction on \(L^2\) by Jensen's inequality and integration. A geometric series
then gives the bounded inverse of the axial block. The other off-diagonal block is bounded because it has
bounded kernel and finite measure. The triangular inverse exists. Reduce \(r\) so that \(\|\mathcal
T_0^{-1}(\mathcal T_x-\mathcal T_0)\|<1/2\). The Neumann series gives a uniformly bounded holomorphic
inverse on that disk.

The part of \(L_0+3\Pi\) with endpoints \(a,b\) equals
\[
 D_x(a)D_x(b)\int K_0(a,t)D_x(t)
     [\mathcal T_x^{-1}(u_x K_0(\cdot,b))](t)\,d\sigma(t).
\]
The two exterior sections have uniform \(L^2\) norms for \emph{each} endpoint by the crossing formula.
The bilinear integral is holomorphic and bounded also when \(a=b\). This gives the version of the kernel
that can be evaluated at a marked cell or on the diagonal. In \(\theta_a\) the factor \(K(a,t)^2\) leaves
\(D_x(t)^2\), which absorbs the single \(1/D_x(t)\) in the measure. In \(s_a\) and in the last sum of
\(\mathcal V\), the fibers and the normalized counting measure on the lines are fixed. The integral
formulas in Appendix~\ref{app:predictor} therefore show that \(\mathcal V\) is holomorphic.

\paragraph{The rook functional.}
At \(x=0\), the scaling of \(v_0=u_Du_A\) is \(\nu_1\) by its factorization. In \(C[0,1]\) the equation
for the logarithm \(z\) of the symmetric factor is \(e^{z(r)}\int v_x(r,q)e^{z(q)}\,dq=1\). Its
derivative at the initial solution is \(I+\mathcal P\), where \(\mathcal P\) has kernel \(\nu_1\). It
acts as multiplication by two on the constants. On the mean-zero subspace, \(\|\mathcal
P\|_\infty\le1-c\) by the same doubly stochastic decomposition. The implicit function theorem in this
complex space gives a holomorphic scaling on a smaller disk. By uniqueness its real branch agrees with
the positive scaling. For the centered kernel \(\centrado=p_x-1\) of Appendix~\ref{app:predictor},
\(\|\centrado(0)\|_{2\to2}\le1-c\). On reducing the disk, \((I-\centrado^2)^{-1}\) is bounded and
holomorphic. The kernel of \(T=\centrado(I-\centrado^2)^{-1}\centrado\) is evaluated between two bounded
sections of \(\centrado\), including on the diagonal, as in the preceding paragraph. Kernel products and
\(\operatorname{tr}(T^2)=\iint T(r,q)T(q,r)\,dr\,dq\) preserve holomorphy. Thus \(F\), and hence \(G\),
are holomorphic near zero.

\paragraph{Maximum order in the contractions.}
In a direct product, the expectation of the indicators is the product of the \(U_\ell\) for the distinct
unmarked lines. It is holomorphic even for two crossings sharing a transversal. This product is used
before any division by the shared availability. Area means use \(\nu_1\,di\), and the fiber of a marked
cell uses \(\eta_\lambda\nu_1\,dy\), both bounded. Each free integral \(d\rho\) costs at most one factor
\(x^{-1}\). The table counts those factors in the physical expressions in Appendix~\ref{app:interior},
including the squares and their two corrections where applicable.
\begin{center}\small
\begin{tabular}{@{}lp{.54\linewidth}@{}}\toprule
Contraction & Possible sources of poles\\ \midrule
\(G_{\rm cof}, f2E, f2F,f3F,f3c,f4,f5c,f5F,f7\) & Only areas, centers, and marked fibers.\\
\(f2z\) & One free line and the bounded multiplier \(u_f/U_f\).\\
\(f6\) & \(\Lambda,\omega\) use crossings and fibers. The response \(q\) integrates over one free line.\\
\(C_{3,\rm sq}\) & The explicit power \(x^{n_{RC}-3}\) costs at most three factors.\\
\(Q_1+\Delta_1\) & One copy of \(\mathcal E_2\) integrates over at most one free line.\\
\(Q_2+\Delta_2\) & Two exterior lines and at most one interior line in each copy of \(\mathcal E_2\).\\
\bottomrule
\end{tabular}
\end{center}
The last two rows use \(L_1,L_2^{\rm int}\) before the products are decomposed into covariances. The
discrete sum with \(\mathcal S(w,u)\) runs over the transversals of the crossings \(u\cap j\), whose
number is bounded. Thus each product in the table, multiplied by \(x^4\), is holomorphic under fixed
measures and has an integrable constant majorant on a smaller disk. Integrating the series by Cauchy's
formula on an interior circle gives the holomorphy of \(x^4d(x)\). Since \(\numerr=d+G-xG'\), the claim
follows.
\end{proof}
